\begin{abstract}
Byzantine fault-tolerant (BFT) consensus secures agreement on transaction logs
but generally does not constrain the order chosen within a valid proposal,
leaving room for transaction-order manipulation and maximal extractable value
(MEV).  Graph-based fair-ordering protocols constrain this discretion by
aggregating replicas' receipt observations, yet their validation often remains
symmetric: followers reconstruct the dependency graph and repeat its
extraction before accepting a proposal.  We present AUTIG, a modular
fair-ordering layer that separates authoritative cumulative evidence from
graph materialization.  Replicas maintain authenticated, prefix-consistent
receipt-order evidence, while only the order leader materializes the resulting
state as a persistent Unconfirmed-Transaction Incremental Graph and
incrementally refreshes affected relations.  Each proposal carries a
structural certificate for a deterministic safe closure.  Followers evaluate
pair relations from local evidence to validate the selected order, its
incoming boundary, and why excluded transactions remain pending, without
materializing the complete active graph or repeating full-graph extraction.  We
prove authenticated-state agreement, $\gamma$-batch-order-fairness, pruning
and recovery safety, and conditional progress under stabilized ancestry.  In
ordering-layer experiments on AWS EC2, AUTIG improves throughput by
1.6$\times$--3.4$\times$ and reduces ordering latency by 28\%--40\% over a
Themis-style baseline in the reported 30--60-replica configurations.
\end{abstract}

\IEEEpeerreviewmaketitle

%-------------------------------------------------------------------------------
\section{Introduction}
\label{sec:introduction}
%-------------------------------------------------------------------------------

Byzantine fault-tolerant (BFT) state-machine replication allows a set of
replicas to agree on a common, totally ordered transaction log despite
Byzantine participants~\cite{schneider1990implementing,xu2023survey,
zhang2024reaching,
castro1999practical,yin2019hotstuff}.  Agreement on a valid proposal, however,
does not generally distinguish among its admissible transaction orders.
Because multiple orders may satisfy validity, a proposal leader retains
discretion over the relative placement of included transactions without
violating consensus safety.  In blockchain and financial applications, this
discretion creates opportunities for transaction-order manipulation and
maximal extractable value (MEV)~\cite{daian2020flash,eskandari2019sok,
gramoli2024resilience,
qin2022quantifying,zhou2021high,baum2022sok,heimbach2022eliminating,
heimbach2024non}.  Order-fairness treats this discretion as an ordering-layer
correctness concern: the committed order should reflect information reported
by multiple replicas rather than the preference of one proposal
leader~\cite{kelkar2020order,orda2021enforcing,kiayias2024ordering,
cachin2022quick,kelkar2022order,zhang2020byzantine,kursawe2020wendy,
li2024sok,chen2024auncel}.

Fair-ordering approaches constrain the output using different signals,
including causal or timing information and relative or weighted receipt
observations~\cite{kelkar2022orderfair,kursawe2020wendy,zhang2020byzantine,
chen2024auncel}.  Among these approaches, graph-based relative-order fairness
provides a concrete mechanism for batch-order guarantees.  Replicas report
signed local precedence observations; threshold-supported pair relations
induce a directed transaction-dependency graph; and strongly connected
components (SCCs) group Condorcet cycles that cannot all be represented by a
strict total order.  This abstraction underlies Aequitas and Themis-style
protocols~\cite{kelkar2020order,kelkar2023themis}: transactions in one SCC may
share a fairness batch, while the condensation graph constrains the order
among different batches.

The graph abstraction also introduces a systems bottleneck.  In representative
graph-based designs, validation remains largely symmetric: followers derive
pair relations from the selected observations, materialize the relevant
dependency graph, and repeat SCC and condensation-graph extraction before
accepting the leader's result.  With $|A|$ active transactions, the relation
space may contain $\Theta(|A|^2)$ pairs, and substantially the same graph work
is repeated across replicas even when a proposal finalizes only part of the
active set.  SpeedyFair separates fair ordering from consensus and overlaps
their execution~\cite{mu2024separation}, but does not remove replica-wide
extraction within the fair-ordering stage.  Improving scalability therefore
requires more than moving ordering work off the consensus path: it requires
breaking the symmetry of graph materialization and extraction without making
the leader's derived graph authoritative.

Making verification asymmetric is not equivalent to trusting the order
leader.  Since fairness evidence accumulates across proposals, a Byzantine
leader must not be able to alter historical pair support, omit an external
predecessor, or hand off an incomplete live state.  Followers therefore need
an authoritative cumulative-evidence state against which the leader's graph
decisions can be checked.  This raises the central question of this work:
\emph{can followers validate a graph-based fair order without persistently
materializing the complete active graph or repeating its full extraction?}

AUTIG answers this question by separating authenticated evidence from its
graph representation.  All replicas retain prefix-consistent receipt-order
evidence from which cumulative pair relations are derived deterministically.
The order leader alone materializes this evidence as a persistent
\emph{Unconfirmed-Transaction Incremental Graph} (UTIG) and refreshes newly
active vertices and affected pairs as evidence arrives.  The UTIG is therefore
a leader-side acceleration structure rather than an independent source of
protocol truth.

For each proposal, AUTIG extracts a deterministic \emph{safe closure}: the
maximal down-closed subset of active transactions whose vertices have reached
the protocol's \emph{Solid} visibility threshold.  The leader attaches a
structural certificate that identifies the selected SCCs and their
deterministic order and supplies a blocker forest for excluded Solid
transactions.  Followers evaluate the required pair predicates from local
evidence to verify the certified SCCs, the complete incoming boundary, and the
blocker witnesses.  They do not persist the complete active dependency graph
or repeat full-graph extraction.  The structural certificate has linear size
in the finalized set and the represented blocker forest, although verification
may still inspect quadratically many pair values in the worst case.

AUTIG integrates this verification path with committed state continuity.  A
locally accepted proposal produces only a candidate state, which becomes
authoritative after the hosting BFT commits the fragment.  Derived graph
caches may be pruned, while authenticated live evidence is retained.
Checkpoints and committed replay preserve this state across recovery and
order-leader handoff.  We prove authenticated-state agreement,
incremental-materialization equivalence, safe-fragment validity,
$\gamma$-batch-order-fairness, pruning safety, and exact recovery.  The
progress result is conditional on stabilized ancestry: it does not provide
unconditional progress when new arrivals indefinitely expand a transaction's
live ancestor set.

We implement AUTIG as an ordering-layer prototype and compare it with a
Themis-style baseline on AWS EC2 with 10--100 replicas.  In the reported
30--60-replica configurations, AUTIG improves throughput by
1.6$\times$--3.4$\times$ and reduces ordering latency by 28\%--40\%.  These
measurements characterize the ordering layer rather than an end-to-end BFT
deployment.  Figure~\ref{fig:autig-overview} summarizes the architectural
distinction.

\begin{figure}[!t]
\centering
\subfloat[]{
  \includegraphics[width=0.8\linewidth]{themis.pdf}
  \label{fig:themis}
}

\vspace{0.4em}

\subfloat[]{
  \includegraphics[width=0.8\linewidth]{autig.pdf}
  \label{fig:autig}
}
\caption{Ordering and validation architectures.  (a) A Themis-style design
reconstructs the proposal graph and extraction state at each replica.
(b) AUTIG replicates authenticated receipt-order evidence at all replicas but
materializes the persistent UTIG only at the order leader.  Followers validate
the selected order through a structural certificate and locally derived pair
relations.}
\label{fig:autig-overview}
\end{figure}

The contributions of this paper are as follows.
\begin{itemize}
  \item We introduce an authenticated cumulative-evidence representation based
  on prefix-consistent receipt logs and unique per-replica first positions.
  It makes pair relations independently derivable at followers while treating
  graph structures as reconstructible, non-authoritative derivatives.

  \item We design a persistent incremental UTIG and prove that refreshing
  newly active vertices and touched pairs is equivalent to rematerializing the
  protocol-defined graph from the authenticated state.

  \item We develop a safe-closure rule and structural certificate that,
  together with locally derived pair relations, let followers verify the
  finalized SCCs, their deterministic order, the complete incoming frontier,
  and blocker witnesses without maintaining the persistent graph.  We prove
  authenticated-state agreement, safe-fragment validity,
  $\gamma$-batch-order-fairness, pruning safety, exact recovery across handoff,
  and conditional progress under stabilized ancestry.

  \item We implement an ordering-layer prototype and evaluate it against a
  Themis-style baseline on AWS EC2 with 10--100 replicas.  In the reported
  30--60-replica configurations, AUTIG improves throughput by
  1.6$\times$--3.4$\times$ and reduces ordering latency by 28\%--40\%.
\end{itemize}

%-------------------------------------------------------------------------------
\section{Background}
\label{sec:background}
%-------------------------------------------------------------------------------

\subsection{BFT Consensus and Transaction Ordering}

A BFT state-machine replication protocol lets $n$ replicas agree on a common
sequence of commands despite up to $f$ Byzantine replicas
~\cite{schneider1990implementing,castro1999practical,yin2019hotstuff}.
Agreement and finality prevent correct replicas from committing conflicting
histories.  After synchrony and under the protocol's leader and availability
assumptions, liveness ensures continued commitment of admissible proposals.
These properties do not determine which valid transaction a leader selects or
where it places that transaction within a proposal.

Consequently, transaction-order manipulation need not violate BFT safety.  A
leader can exploit information about pending transactions to change
application-visible outcomes while proposing an otherwise valid
block~\cite{daian2020flash,heimbach2022sok,malkhi2022maximal}.  Fair ordering adds
constraints on this discretionary ordering step; the hosting BFT remains
responsible for agreement, availability, and finality of the resulting
ordering records.

\subsection{Relative Order Fairness}
\label{sec:background_order_fairness}

Fair-ordering mechanisms differ in the observations used to constrain the
output.  Timestamp-based approaches use clocks or timestamp attestations and
must account for clock skew, network uncertainty, and adversarial timestamps
~\cite{zhang2020byzantine,kursawe2020wendy}.  Relative-order approaches use
only whether one transaction was received before another at each replica
~\cite{kelkar2020order,kelkar2023themis}.  The latter avoids a shared-clock
assumption but must aggregate potentially conflicting local orders.

A strict global order cannot, in general, satisfy every threshold-supported
pair relation.  Even when each replica reports a transitive local sequence,
pairwise aggregation can produce a Condorcet cycle: a threshold supports
$x\prec y$, another supports $y\prec z$, and a third supports $z\prec x$
~\cite{gehrlein1983condorcet}.  Batch-order-fairness resolves this conflict by
allowing cyclically related transactions to share a fairness batch.  If the
required fraction of correct replicas receives $x$ before $y$, then $x$ must
be assigned to the same or an earlier batch than $y$.
Section~\ref{sec:fairness-objective} gives AUTIG's precise Themis-style
$\gamma$-parameterization and integer feasibility condition.

\subsection{Graph-Based Aggregation and Verification}
\label{sec:background_architectures}

Graph-based protocols represent transactions as vertices and
threshold-supported pair preferences as directed edges.  SCCs identify cyclic
preferences and become fairness batches; contracting them yields a DAG whose
topological order constrains different batches.  Aequitas established this
general construction but can suffer weak progress when newly arriving
transactions continue to extend dependency chains~\cite{kelkar2020order}.
Themis develops deferred-ordering machinery for stronger progress
~\cite{kelkar2023themis}.  AUTIG adopts this relative-order fairness objective
and graph semantics while changing the validation architecture; its
safe-closure progress condition is stated separately in
Proposition~\ref{prop:stabilized-progress}.

The standard validation strategy reconstructs the graph from selected signed
observations and repeats the leader's SCC, condensation, and extraction logic.
This symmetry provides a direct validity check but replicates graph
materialization and extraction across the system.  SpeedyFair separates fair
ordering from the main BFT path and overlaps the two activities, improving
scheduling while retaining replica-wide fair-order computation
~\cite{mu2024separation}.  AUTIG addresses a different boundary: followers
retain authenticated receipt evidence but not the persistent graph.  The
order leader incrementally materializes that graph, and followers validate its
selected safe closure through locally evaluated pair predicates and a
structural certificate.
